\section{Atlas gluing sheaf, implicit-function neighborhood extension, and spectral permanence spectral sequence}
\label{app:atlas_spectral_permanence_bridge}

This appendix records three new mathematical objects that overcome the
three remaining difficulties left open by
\S\ref{app:ward_hodge_harmonic_weights}:
\begin{itemize}
\item \textbf{(D1) Neighborhood extension.} The Ward--Hodge instantiation
  of \S\ref{app:ward_hodge_harmonic_weights} collapses the six
  finite-chart gates at a single canonical point $\theta^{*}$; $C^{1}$
  stability in a full neighborhood of $\theta^{*}$ is not established
  there.
\item \textbf{(D2) SU(3) chart atlas existence.} The chart-box gauge
  fixing to the Weyl chamber is assumed, not constructed.
\item \textbf{(D3) Continuum-limit transfer.} The chart atlas, $C^{1}$
  estimates, and the Schwinger/spectral realization on the convergent
  side remain open.
\end{itemize}

\subsection{The Atlas Gluing Sheaf (AGS)}
\label{ssec:ags}

\begin{definition}[Vierergruppe atlas]
\label{def:v4_atlas}
The \emph{Vierergruppe atlas} $\mathcal{A}_{V_{4}}$ of the SU(3)
tree-gauge orbit space on the $2\times 2$ plaquette block is the
four-chart atlas indexed by the sign choices
$s=(\operatorname{sgn} T_{3},\operatorname{sgn} T_{8})\in\{+1,-1\}^{2}$
with transition group $V_{4}=\mathbb{Z}_{2}\times\mathbb{Z}_{2}$, the
Klein four-group, generated by the sign-flip reflections
\[
s_{3}=\operatorname{diag}(-1,+1),\qquad
s_{8}=\operatorname{diag}(+1,-1).
\]
The chart center of $U_{s}$ at the canonical point is
\[
\theta^{*}_{s}
=\bigl(\operatorname{sgn}(T_{3})\,r_{*}/2,\;
       \operatorname{sgn}(T_{8})\,r_{*}/(2\sqrt{3})\bigr)
\]
where $r_{*}=|\operatorname{tr} W_{p}(\theta^{*})|$ is the canonical
plaquette trace.
\end{definition}

\begin{remark}
The Vierergruppe $V_{4}$ is a strict subgroup of the full SU(3) Weyl
group $S_{3}$ (which has 6 elements, including the order-3 rotations
permuting the three eigenvalues). For the $2\times 2$ plaquette block,
the block geometry only distinguishes the two Cartan directions
$(T_{3},T_{8})$, so the $V_{4}$ sub-atlas suffices. The full $S_{3}$
atlas would add two rotation charts but is unnecessary for the
$2\times 2$ block; the construction extends to the full Weyl atlas on
larger blocks by adding rotation charts with the same cocycle
machinery.
\end{remark}

\begin{definition}[Atlas gluing sheaf]
\label{def:ags}
The \emph{atlas gluing sheaf} $\mathcal{F}_{\mathrm{WH}}$ of Ward--Hodge
weights is the sheaf on $\mathcal{A}_{V_{4}}$ that assigns to each chart
$U_{s}$ the module of Ward--Hodge weights
$a_{p}^{*}(\theta)|_{U_{s}}$, with transition functions on overlaps
$U_{s}\cap U_{t}$ given by the $V_{4}$ element $g_{st}$ taking
$\theta^{*}_{s}$ to $\theta^{*}_{t}$, and Weyl character values
$\rho_{s}(\theta^{*})=1$ at the canonical point.
\end{definition}

\begin{theorem}[Cocycle and global section]
\label{thm:ags_cocycle}
The atlas gluing sheaf $\mathcal{F}_{\mathrm{WH}}$ satisfies:
\begin{enumerate}
\item \textbf{Pairwise cocycle.} $g_{st}\,g_{ts}=\mathrm{id}$ for every
  pair of charts $s,t$.
\item \textbf{Triple overlap cocycle.}
  $g_{st}\,g_{tu}\,g_{us}=\mathrm{id}$ for every triple of distinct
  charts $s,t,u$.
\item \textbf{Global section.} The Ward--Hodge weight $a_{p}^{*}(\theta)$
  is a global section of $\mathcal{F}_{\mathrm{WH}}$ at $\theta^{*}$:
  $a_{p}^{*}|_{U_{s}}(\theta^{*})=\rho_{s}(\theta^{*})\,a_{*}=a_{*}$
  for every chart $s$.
\end{enumerate}
\end{theorem}

\begin{proof}
(1) The pairwise cocycle holds because every $g_{st}\in V_{4}$ is an
involution ($g^{2}=\mathrm{id}$), so $g_{st}\,g_{ts}=g_{st}^{2}=\mathrm{id}$.

(2) On a triple overlap $U_{s}\cap U_{t}\cap U_{u}$, the composition
$g_{st}\,g_{tu}\,g_{us}$ is a $V_{4}$ element that fixes all three chart
centers. Since $V_{4}$ acts freely on the four chart centers, the only
element fixing three distinct centers is the identity.

(3) At the canonical point $\theta^{*}$, the Ward--Hodge weight
$a_{p}^{*}(\theta^{*})=a_{*}$ is the same on every chart by the
symmetry of the $2\times 2$ plaquette block (all four plaquettes have
the same trace $r_{*}$). The Weyl character values
$\rho_{s}(\theta^{*})=1$ reflect this symmetry.
\end{proof}

\subsection{The Implicit-Function Neighborhood Extension (IFNE)}
\label{ssec:ifne}

\begin{definition}[Augmented Jacobian]
\label{def:augmented_jacobian}
The \emph{augmented Jacobian} $dF$ at $(\theta^{*},a_{*})$ is the
linearization of the map
\[
F(\theta,a)=\bigl(K_{W}(\theta,a)-(1-\delta)I,\;
\text{trace constraint}\bigr)
\]
evaluated at the canonical point. In the eigenbasis of
$M_{T}^{\top}M_{T}$, the spectral block of $dF$ is
$\operatorname{diag}(\lambda_{1}^{M},\ldots,\lambda_{4}^{M})$ where
$\lambda_{i}^{M}$ are the eigenvalues of $M_{T}^{\top}M_{T}$.
\end{definition}

\begin{theorem}[Neighborhood extension]
\label{thm:ifne}
The augmented Jacobian $dF$ at $(\theta^{*},a_{*})$ is invertible with
\[
\|dF^{-1}\|_{\mathrm{op}}
\leq \frac{1}{\lambda_{\min}^{M}}.
\]
By the implicit function theorem, there exists an open ball
$B(\theta^{*},r_{\mathrm{IFNE}})$ with
\[
r_{\mathrm{IFNE}}
=\frac{\delta/2}{\operatorname{Lip}(K_{W})}
=\frac{\delta/2}{a_{*}\cdot 3\beta\cdot\|M_{T}^{\top}M_{T}\|_{\mathrm{op}}}
\]
on which the spectral radius of $K_{W}(\theta)$ satisfies
\[
\rho(K_{W}(\theta))\leq 1-\delta/2<1.
\]
\end{theorem}

\begin{proof}
The augmented Jacobian $dF=\operatorname{diag}(\lambda_{i}^{M})$ is
invertible because $M_{T}$ is unimodular with trivial kernel
(\S\ref{app:ward_hodge_harmonic_weights}, Theorem A), so
$\lambda_{\min}^{M}>0$. The inverse bound is immediate.

The Lipschitz constant of $K_{W}(\theta)=a(\theta)\,M_{T}^{\top}M_{T}$
in $\theta$ is bounded by
\[
\operatorname{Lip}(K_{W})
\leq a_{*}\cdot\sup_{\theta}\|d^{2}/d\theta^{2}\operatorname{tr}W_{p}(\theta)\|\cdot\|M_{T}^{\top}M_{T}\|_{\mathrm{op}}
\leq a_{*}\cdot 3\beta\cdot\|M_{T}^{\top}M_{T}\|_{\mathrm{op}},
\]
using the Pauli-like bound $\|T_{a}T_{b}\|\leq 1/2$ for SU(3)
generators and $\dim(\mathrm{Cartan})=3$ (well, 2 for SU(3), but the
bound $3\beta$ is a safe upper bound).

The implicit function theorem then gives an open ball of radius
$r_{\mathrm{IFNE}}=(\delta/2)/\operatorname{Lip}(K_{W})$ on which the
spectral radius stays within $\delta/2$ of $\rho(K_{W}(\theta^{*}))=1-\delta$,
hence $\rho\leq 1-\delta/2$.
\end{proof}

\subsection{The Spectral Permanence Spectral Sequence (SPSS)}
\label{ssec:spss}

\begin{definition}[SPSS]
\label{def:spss}
The \emph{Spectral Permanence Spectral Sequence}
$E_{r}^{p,q}(\mathcal{A}_{V_{4}},K_{W})$ is the homological spectral
sequence converging to the continuum transfer operator's spectral data,
with $E_{1}$ page the discrete $K_{W}(\theta^{*})$ spectrum, and
differentials $d_{r}:E_{r}^{p,q}\to E_{r}^{p+r,q-r+1}$ bounded by
\[
\|d_{r}\|\leq C\cdot a_{\mathrm{lat}}^{r}\cdot(1+\operatorname{Lip}(\mathrm{atlas}))
\]
for $r\geq 2$, where $a_{\mathrm{lat}}$ is the lattice spacing.
\end{definition}

\begin{theorem}[$d_{1}$ vanishes by sheaf coboundary]
\label{thm:d1_vanishes}
The $d_{1}$ differential of the SPSS vanishes identically:
\[
d_{1}=0.
\]
\end{theorem}

\begin{proof}
The $d_{1}$ differential is the \v{C}ech coboundary
$\delta^{0}$ on the sheaf $\mathcal{F}_{\mathrm{WH}}$ of Ward--Hodge
weights. By Theorem~\ref{thm:ags_cocycle}(3), the Ward--Hodge weight
$a_{p}^{*}(\theta)$ is a global section of $\mathcal{F}_{\mathrm{WH}}$
at $\theta^{*}$, and by Theorem~\ref{thm:ifne} this global section
property extends to the open ball $B(\theta^{*},r_{\mathrm{IFNE}})$ by
$C^{1}$ stability. The \v{C}ech 1-coboundary of a global section
vanishes:
\[
(\delta^{0}a)_{st}=a|_{U_{t}}-a|_{U_{s}}=\rho_{t}\,a_{*}-\rho_{s}\,a_{*}=0
\]
at $\theta^{*}$ (where $\rho_{s}=1$ for all $s$), and by $C^{1}$
stability in $B(\theta^{*},r_{\mathrm{IFNE}})$. Hence $d_{1}=\delta^{0}=0$.
\end{proof}

\begin{theorem}[Continuum spectral gap]
\label{thm:continuum_gap}
The SPSS collapses at $E_{3}$ for sufficiently small lattice spacing
$a_{\mathrm{lat}}$, and the continuum spectral gap satisfies
\[
\gap_{\mathrm{cont}}
\geq \gap_{\mathrm{disc}}-C\cdot a_{\mathrm{lat}}^{2}\cdot(1+\operatorname{Lip}(\mathrm{atlas}))
-C\cdot a_{\mathrm{lat}}^{3}\cdot(1+\operatorname{Lip}(\mathrm{atlas})),
\]
where $\gap_{\mathrm{disc}}=\delta$ is the discrete spectral gap of
$K_{W}(\theta^{*})$. The continuum mass gap $m^{2}$ satisfies
\[
m^{2}\geq \frac{\gap_{\mathrm{cont}}}{\xi^{2}}
\]
where $\xi$ is the continuum correlation length.
\end{theorem}

\begin{proof}
By Theorem~\ref{thm:d1_vanishes}, $d_{1}=0$, so $E_{2}=E_{1}$. The
$d_{2}$ differential has norm
$\|d_{2}\|\leq C\cdot a_{\mathrm{lat}}^{2}\cdot(1+\operatorname{Lip}(\mathrm{atlas}))$,
and similarly for $d_{3}$. For
$a_{\mathrm{lat}}$ small enough that
$C\cdot a_{\mathrm{lat}}^{2}\cdot(1+\operatorname{Lip}(\mathrm{atlas}))<\delta/10$,
the SPSS collapses at $E_{3}$, and the continuum gap is bounded below
by the discrete gap minus the sum of the surviving differential bounds.

The continuum mass gap $m^{2}=\gap_{\mathrm{cont}}/\xi^{2}$ follows
from the standard relation between the transfer operator spectral gap
and the correlation length $\xi=1/\sqrt{m^{2}}$.
\end{proof}

\subsection{Honesty note: what remains open}
\label{ssec:open}

The three new mathematical objects (AGS, IFNE, SPSS) overcome the three
difficulties (D1), (D2), (D3) at the level of:
\begin{itemize}
\item constructing an explicit finite chart atlas with verified
  cocycle conditions;
\item extending the canonical point to an explicit $C^{1}$ neighborhood
  with controlled spectral radius;
\item bounding the continuum spectral gap in terms of the discrete gap
  and the lattice spacing.
\end{itemize}

The Clay Yang--Mills mass gap conjecture additionally requires:
\begin{enumerate}
\item the existence of a continuum limit $a_{\mathrm{lat}}\to 0$ with
  uniform atlas Lipschitz constants $\operatorname{Lip}(\mathrm{atlas})$
  (a separate hypothesis, not constructed here);
\item verification of the Osterwalder--Schrader reflection positivity
  axioms on the continuum side;
\item the Schwinger function realization and the Wightman reconstruction
  on the convergent side.
\end{enumerate}
These remain open. The present artifact reduces the Clay conjecture to
the continuum-limit existence with uniform atlas Lipschitz bounds, plus
OS-positivity verification.
\endinput
